%% kapa: exam=1 round=33 year=2022 period=1 form=A subject=civil_law subject_ko=민법 seq=9 number=9
%% source: Q-Net 감정평가사 기출문제 (cst00309) artlSeq=5209997 file=raw/1008/5209997/2022년도 제33회 감정평가사 1차시험 1교시(민법,경제학,부동산학).pdf
%% conversion: kordoc 4.24.51 -> tools/gen1.py (auto); review=proofread-sonnet; flags=none
\begin{question}{9}
불공정한 법률행위에 관한 설명으로 옳지 않은 것은? (다툼이 있으면 판례에 따름)
\begin{choices}
  \item 급부와 반대급부 사이의 현저한 불균형은 그 무효를 주장하는 자가 증명해야 한다.
  \item 무경험은 어느 특정 영역에서의 경험부족이 아니라 거래일반에 대한 경험부족을 의미한다.
  \item 대리인에 의한 법률행위의 경우, 궁박 상태에 있었는지 여부는 본인을 기준으로 판단한다.
  \item 불공정한 법률행위로서 무효인 경우, 원칙적으로 추인에 의하여 유효로 될 수 없다.
  \item 경매절차에서 매각대금이 시가보다 현저히 저렴한 경우, 그 경매는 불공정한 법률행위로서 무효이다.
\end{choices}
\end{question}
\answer{5}
